\IfFileExists{stacks-project.cls}{%
\documentclass{stacks-project}
}{%
\documentclass{amsart}
}

% For dealing with references we use the comment environment
\usepackage{verbatim}
\newenvironment{reference}{\comment}{\endcomment}
%\newenvironment{reference}{}{}
\newenvironment{slogan}{\comment}{\endcomment}
\newenvironment{history}{\comment}{\endcomment}

% For commutative diagrams we use Xy-pic
\usepackage[all]{xy}

% We use 2cell for 2-commutative diagrams.
\xyoption{2cell}
\UseAllTwocells

% We use multicol for the list of chapters between chapters
\usepackage{multicol}

% This is generally recommended for better output
\usepackage{lmodern}
\usepackage[T1]{fontenc}

% For cross-file-references
\usepackage{xr-hyper}

% Package for hypertext links:
\usepackage{hyperref}

% For any local file, say "hello.tex" you want to link to please
% use \externaldocument[hello-]{hello}
\externaldocument[introduction-]{introduction}
\externaldocument[conventions-]{conventions}
\externaldocument[sets-]{sets}
\externaldocument[categories-]{categories}
\externaldocument[topology-]{topology}
\externaldocument[sheaves-]{sheaves}
\externaldocument[sites-]{sites}
\externaldocument[stacks-]{stacks}
\externaldocument[fields-]{fields}
\externaldocument[algebra-]{algebra}
\externaldocument[brauer-]{brauer}
\externaldocument[homology-]{homology}
\externaldocument[derived-]{derived}
\externaldocument[simplicial-]{simplicial}
\externaldocument[more-algebra-]{more-algebra}
\externaldocument[smoothing-]{smoothing}
\externaldocument[modules-]{modules}
\externaldocument[sites-modules-]{sites-modules}
\externaldocument[injectives-]{injectives}
\externaldocument[cohomology-]{cohomology}
\externaldocument[sites-cohomology-]{sites-cohomology}
\externaldocument[dga-]{dga}
\externaldocument[dpa-]{dpa}
\externaldocument[sdga-]{sdga}
\externaldocument[hypercovering-]{hypercovering}
\externaldocument[schemes-]{schemes}
\externaldocument[constructions-]{constructions}
\externaldocument[properties-]{properties}
\externaldocument[morphisms-]{morphisms}
\externaldocument[coherent-]{coherent}
\externaldocument[divisors-]{divisors}
\externaldocument[limits-]{limits}
\externaldocument[varieties-]{varieties}
\externaldocument[topologies-]{topologies}
\externaldocument[descent-]{descent}
\externaldocument[perfect-]{perfect}
\externaldocument[more-morphisms-]{more-morphisms}
\externaldocument[flat-]{flat}
\externaldocument[groupoids-]{groupoids}
\externaldocument[more-groupoids-]{more-groupoids}
\externaldocument[etale-]{etale}
\externaldocument[chow-]{chow}
\externaldocument[intersection-]{intersection}
\externaldocument[pic-]{pic}
\externaldocument[weil-]{weil}
\externaldocument[adequate-]{adequate}
\externaldocument[dualizing-]{dualizing}
\externaldocument[duality-]{duality}
\externaldocument[discriminant-]{discriminant}
\externaldocument[derham-]{derham}
\externaldocument[local-cohomology-]{local-cohomology}
\externaldocument[algebraization-]{algebraization}
\externaldocument[curves-]{curves}
\externaldocument[resolve-]{resolve}
\externaldocument[models-]{models}
\externaldocument[functors-]{functors}
\externaldocument[equiv-]{equiv}
\externaldocument[pione-]{pione}
\externaldocument[etale-cohomology-]{etale-cohomology}
\externaldocument[proetale-]{proetale}
\externaldocument[relative-cycles-]{relative-cycles}
\externaldocument[more-etale-]{more-etale}
\externaldocument[trace-]{trace}
\externaldocument[crystalline-]{crystalline}
\externaldocument[spaces-]{spaces}
\externaldocument[spaces-properties-]{spaces-properties}
\externaldocument[spaces-morphisms-]{spaces-morphisms}
\externaldocument[decent-spaces-]{decent-spaces}
\externaldocument[spaces-cohomology-]{spaces-cohomology}
\externaldocument[spaces-limits-]{spaces-limits}
\externaldocument[spaces-divisors-]{spaces-divisors}
\externaldocument[spaces-over-fields-]{spaces-over-fields}
\externaldocument[spaces-topologies-]{spaces-topologies}
\externaldocument[spaces-descent-]{spaces-descent}
\externaldocument[spaces-perfect-]{spaces-perfect}
\externaldocument[spaces-more-morphisms-]{spaces-more-morphisms}
\externaldocument[spaces-flat-]{spaces-flat}
\externaldocument[spaces-groupoids-]{spaces-groupoids}
\externaldocument[spaces-more-groupoids-]{spaces-more-groupoids}
\externaldocument[bootstrap-]{bootstrap}
\externaldocument[spaces-pushouts-]{spaces-pushouts}
\externaldocument[spaces-chow-]{spaces-chow}
\externaldocument[groupoids-quotients-]{groupoids-quotients}
\externaldocument[spaces-more-cohomology-]{spaces-more-cohomology}
\externaldocument[spaces-simplicial-]{spaces-simplicial}
\externaldocument[spaces-duality-]{spaces-duality}
\externaldocument[formal-spaces-]{formal-spaces}
\externaldocument[restricted-]{restricted}
\externaldocument[spaces-resolve-]{spaces-resolve}
\externaldocument[formal-defos-]{formal-defos}
\externaldocument[defos-]{defos}
\externaldocument[cotangent-]{cotangent}
\externaldocument[examples-defos-]{examples-defos}
\externaldocument[algebraic-]{algebraic}
\externaldocument[examples-stacks-]{examples-stacks}
\externaldocument[stacks-sheaves-]{stacks-sheaves}
\externaldocument[criteria-]{criteria}
\externaldocument[artin-]{artin}
\externaldocument[quot-]{quot}
\externaldocument[stacks-properties-]{stacks-properties}
\externaldocument[stacks-morphisms-]{stacks-morphisms}
\externaldocument[stacks-limits-]{stacks-limits}
\externaldocument[stacks-cohomology-]{stacks-cohomology}
\externaldocument[stacks-perfect-]{stacks-perfect}
\externaldocument[stacks-introduction-]{stacks-introduction}
\externaldocument[stacks-more-morphisms-]{stacks-more-morphisms}
\externaldocument[stacks-geometry-]{stacks-geometry}
\externaldocument[moduli-]{moduli}
\externaldocument[moduli-curves-]{moduli-curves}
\externaldocument[examples-]{examples}
\externaldocument[exercises-]{exercises}
\externaldocument[guide-]{guide}
\externaldocument[desirables-]{desirables}
\externaldocument[coding-]{coding}
\externaldocument[obsolete-]{obsolete}
\externaldocument[fdl-]{fdl}
\externaldocument[index-]{index}

% Theorem environments.
%
\theoremstyle{plain}
\newtheorem{theorem}[subsection]{Theorem}
\newtheorem{proposition}[subsection]{Proposition}
\newtheorem{lemma}[subsection]{Lemma}

\theoremstyle{definition}
\newtheorem{definition}[subsection]{Definition}
\newtheorem{example}[subsection]{Example}
\newtheorem{exercise}[subsection]{Exercise}
\newtheorem{situation}[subsection]{Situation}

\theoremstyle{remark}
\newtheorem{remark}[subsection]{Remark}
\newtheorem{remarks}[subsection]{Remarks}

\numberwithin{equation}{subsection}

% Macros
%
\def\lim{\mathop{\mathrm{lim}}\nolimits}
\def\colim{\mathop{\mathrm{colim}}\nolimits}
\def\Spec{\mathop{\mathrm{Spec}}}
\def\Hom{\mathop{\mathrm{Hom}}\nolimits}
\def\Ext{\mathop{\mathrm{Ext}}\nolimits}
\def\SheafHom{\mathop{\mathcal{H}\!\mathit{om}}\nolimits}
\def\SheafExt{\mathop{\mathcal{E}\!\mathit{xt}}\nolimits}
\def\Sch{\mathit{Sch}}
\def\Mor{\mathop{\mathrm{Mor}}\nolimits}
\def\Ob{\mathop{\mathrm{Ob}}\nolimits}
\def\Sh{\mathop{\mathit{Sh}}\nolimits}
\def\NL{\mathop{N\!L}\nolimits}
\def\CH{\mathop{\mathrm{CH}}\nolimits}
\def\proetale{{pro\text{-}\acute{e}tale}}
\def\etale{{\acute{e}tale}}
\def\QCoh{\mathit{QCoh}}
\def\Ker{\mathop{\mathrm{Ker}}}
\def\Im{\mathop{\mathrm{Im}}}
\def\Coker{\mathop{\mathrm{Coker}}}
\def\Coim{\mathop{\mathrm{Coim}}}

% Boxtimes
%
\DeclareMathSymbol{\boxtimes}{\mathbin}{AMSa}{"02}

%
% Macros for moduli stacks/spaces
%
\def\QCohstack{\mathcal{QC}\!\mathit{oh}}
\def\Cohstack{\mathcal{C}\!\mathit{oh}}
\def\Spacesstack{\mathcal{S}\!\mathit{paces}}
\def\Quotfunctor{\mathrm{Quot}}
\def\Hilbfunctor{\mathrm{Hilb}}
\def\Curvesstack{\mathcal{C}\!\mathit{urves}}
\def\Polarizedstack{\mathcal{P}\!\mathit{olarized}}
\def\Complexesstack{\mathcal{C}\!\mathit{omplexes}}
% \Pic is the operator that assigns to X its picard group, usage \Pic(X)
% \Picardstack_{X/B} denotes the Picard stack of X over B
% \Picardfunctor_{X/B} denotes the Picard functor of X over B
\def\Pic{\mathop{\mathrm{Pic}}\nolimits}
\def\Picardstack{\mathcal{P}\!\mathit{ic}}
\def\Picardfunctor{\mathrm{Pic}}
\def\Deformationcategory{\mathcal{D}\!\mathit{ef}}
\usepackage{cleveref}
\crefname{originaldefinition}{original definition}{original definitions}
\Crefname{originaldefinition}{Original definition}{Original definitions}
\makeatletter
\newlabel{situation-complexes}{{0DLE}{source}{}{sourcepointer-1}{}}
\newlabel{situation-complexes@cref}{{[originaldefinition][0][]0DLE}{[1][0][]source}{}{}{}}
\makeatother
\crefname{originallemma}{original lemma}{original lemmas}
\Crefname{originallemma}{Original lemma}{Original lemmas}
\makeatletter
\newlabel{more-morphisms-lemma-pushout-along-thickening}{{07RT}{source}{}{sourcepointer-2}{}}
\newlabel{more-morphisms-lemma-pushout-along-thickening@cref}{{[originallemma][0][]07RT}{[1][0][]source}{}{}{}}
\makeatother
\crefname{originallemma}{original lemma}{original lemmas}
\Crefname{originallemma}{Original lemma}{Original lemmas}
\makeatletter
\newlabel{spaces-pushouts-lemma-pushout-along-thickening-schemes}{{07SY}{source}{}{sourcepointer-3}{}}
\newlabel{spaces-pushouts-lemma-pushout-along-thickening-schemes@cref}{{[originallemma][0][]07SY}{[1][0][]source}{}{}{}}
\makeatother
\crefname{originalremark}{original remark}{original remarks}
\Crefname{originalremark}{Original remark}{Original remarks}
\makeatletter
\newlabel{remark-complexes-base-change}{{0DLI}{source}{}{sourcepointer-4}{}}
\newlabel{remark-complexes-base-change@cref}{{[originalremark][0][]0DLI}{[1][0][]source}{}{}{}}
\makeatother
\crefname{originallemma}{original lemma}{original lemmas}
\Crefname{originallemma}{Original lemma}{Original lemmas}
\makeatletter
\newlabel{artin-lemma-pushout}{{07WN}{source}{}{sourcepointer-5}{}}
\newlabel{artin-lemma-pushout@cref}{{[originallemma][0][]07WN}{[1][0][]source}{}{}{}}
\makeatother
\crefname{originallemma}{original lemma}{original lemmas}
\Crefname{originallemma}{Original lemma}{Original lemmas}
\makeatletter
\newlabel{spaces-pushouts-lemma-equivalence-categories-spaces-pushout-flat}{{07W3}{source}{}{sourcepointer-6}{}}
\newlabel{spaces-pushouts-lemma-equivalence-categories-spaces-pushout-flat@cref}{{[originallemma][0][]07W3}{[1][0][]source}{}{}{}}
\makeatother
\crefname{originallemma}{original lemma}{original lemmas}
\Crefname{originallemma}{Original lemma}{Original lemmas}
\makeatletter
\newlabel{spaces-pushouts-lemma-pushout-along-thickening-derived}{{0DL7}{source}{}{sourcepointer-7}{}}
\newlabel{spaces-pushouts-lemma-pushout-along-thickening-derived@cref}{{[originallemma][0][]0DL7}{[1][0][]source}{}{}{}}
\makeatother
\crefname{originallemma}{original lemma}{original lemmas}
\Crefname{originallemma}{Original lemma}{Original lemmas}
\makeatletter
\newlabel{lemma-complexes}{{0DLD}{source}{}{sourcepointer-8}{}}
\newlabel{lemma-complexes@cref}{{[originallemma][0][]0DLD}{[1][0][]source}{}{}{}}
\makeatother
\crefname{originallemma}{original lemma}{original lemmas}
\Crefname{originallemma}{Original lemma}{Original lemmas}
\makeatletter
\newlabel{spaces-more-morphisms-lemma-thickening-relatively-perfect}{{0DL4}{source}{}{sourcepointer-9}{}}
\newlabel{spaces-more-morphisms-lemma-thickening-relatively-perfect@cref}{{[originallemma][0][]0DL4}{[1][0][]source}{}{}{}}
\makeatother
\title{RS* descent for the complexes stack along pushouts along thickenings}
\author{The Stacks Project Authors}
\date{stacks-project:quot.tex:local-snapshot-20261012; GNU FDL 1.2 or later}
\begin{document}
\maketitle
\setcounter{subsection}{1}
\noindent\textit{Editorial note.} All delivered mathematical text is the author's own, copied from the pinned Stacks Project chapter file quot.tex; no statement or proof step was rewritten, reordered or removed. The counted span is the author environment and its own proof environment. References inside the span are handed to the assembler as navigation stubs that carry the canonical tag of the cited tag only; whether each of them is load-bearing is recorded on the stub itself.\par\medskip
\noindent\textit{Allowed standard prerequisites (editorial).}\begin{itemize}
\item Situation situation-complexes (canonical tag 0DLE). cited at statement line 5666. Standing situation.
\item Lemma more-morphisms-lemma-pushout-along-thickening (canonical tag 07RT). cited at statement lines 5674-5675. The pushout square exists.
\item Lemma spaces-pushouts-lemma-pushout-along-thickening-schemes (canonical tag 07SY). cited at proof lines 5692-5693 and 5723-5724. Base bijection; alternative route for bijectivity.
\item Remark remark-complexes-base-change (canonical tag 0DLI). cited at proof lines 5705-5706. Base change reduction to B = Y'.
\item Lemma artin-lemma-pushout (canonical tag 07WN). cited at proof lines 5721-5722. Bijectivity of Isom over the pushout for algebraic spaces.
\item Lemma spaces-pushouts-lemma-equivalence-categories-spaces-pushout-flat (canonical tag 07W3). cited at proof lines 5747-5748. X\_Z, X\_Y, X\_\{Z'\}, X\_\{Y'\} form a pushout square.
\item Lemma spaces-pushouts-lemma-pushout-along-thickening-derived (canonical tag 0DL7). cited at proof lines 5749-5750. Descent of the derived object E\_\{Y'\}.
\item Lemma lemma-complexes (canonical tag 0DLD). cited at proof lines 5756-5757. Condition (2) is a property on points.
\item Lemma spaces-more-morphisms-lemma-thickening-relatively-perfect (canonical tag 0DL4). cited at proof lines 5759-5760. E\_\{Y'\} is Y'-perfect.
\end{itemize}\medskip
\ifdefined\hypertarget\hypertarget{sourcepointer-1}{}\fi
\noindent\textit{Source reference.} Original definition 0DLE: sources/stacks-chow/tags.txt, source lines 17044--17044. Cited at statement line 5666: the lemma is phrased in this situation; load-bearing.\par\medskip
\ifdefined\hypertarget\hypertarget{sourcepointer-2}{}\fi
\noindent\textit{Source reference.} Original lemma 07RT: sources/stacks-chow/tags.txt, source lines 9878--9878. Cited at statement lines 5674-5675: existence of the pushout square; load-bearing.\par\medskip
\ifdefined\hypertarget\hypertarget{sourcepointer-3}{}\fi
\noindent\textit{Source reference.} Original lemma 07SY: sources/stacks-chow/tags.txt, source lines 9918--9918. Cited at proof lines 5692-5693 (base bijection) and 5723-5724 (alternative for bijectivity); load-bearing twice.\par\medskip
\ifdefined\hypertarget\hypertarget{sourcepointer-4}{}\fi
\noindent\textit{Source reference.} Original remark 0DLI: sources/stacks-chow/tags.txt, source lines 17048--17048. Cited at proof lines 5705-5706: replace B by Y' and X by X x\_B Y'; load-bearing for the reduction.\par\medskip
\ifdefined\hypertarget\hypertarget{sourcepointer-5}{}\fi
\noindent\textit{Source reference.} Original lemma 07WN: sources/stacks-chow/tags.txt, source lines 10048--10048. Cited at proof lines 5721-5722: bijectivity of the Isom map for algebraic spaces; load-bearing for full faithfulness.\par\medskip
\ifdefined\hypertarget\hypertarget{sourcepointer-6}{}\fi
\noindent\textit{Source reference.} Original lemma 07W3: sources/stacks-chow/tags.txt, source lines 10028--10028. Cited at proof lines 5747-5748: X\_Z -> X\_Y affine and X\_Z -> X\_\{Z'\} thickening pushout; load-bearing.\par\medskip
\ifdefined\hypertarget\hypertarget{sourcepointer-7}{}\fi
\noindent\textit{Source reference.} Original lemma 0DL7: sources/stacks-chow/tags.txt, source lines 17037--17037. Cited at proof lines 5749-5750: existence of E\_\{Y'\} with compatible isomorphisms; load-bearing.\par\medskip
\ifdefined\hypertarget\hypertarget{sourcepointer-8}{}\fi
\noindent\textit{Source reference.} Original lemma 0DLD: sources/stacks-chow/tags.txt, source lines 17043--17043. Cited at proof lines 5756-5757: condition (2) is a property on points; load-bearing.\par\medskip
\ifdefined\hypertarget\hypertarget{sourcepointer-9}{}\fi
\noindent\textit{Source reference.} Original lemma 0DL4: sources/stacks-chow/tags.txt, source lines 17034--17034. Cited at proof lines 5759-5760: E\_\{Y'\} is Y'-perfect; load-bearing.\par\medskip
\begin{lemma}
\label{lemma-complexes-RS-star}
In Situation \ref{situation-complexes}. Let
$$
\xymatrix{
Z \ar[r] \ar[d] & Z' \ar[d] \\
Y \ar[r] & Y'
}
$$
be a pushout in the category of schemes over $S$ where
$Z \to Z'$ is a finite order thickening and $Z \to Y$ is affine, see
More on Morphisms, Lemma \ref{more-morphisms-lemma-pushout-along-thickening}.
Then the functor on fibre categories
$$
\Complexesstack_{X/B, Y'}
\longrightarrow
\Complexesstack_{X/B, Y}
\times_{\Complexesstack_{X/B, Z}}
\Complexesstack_{X/B, Z'}
$$
is an equivalence.
\end{lemma}

\begin{proof}
Observe that the corresponding map
$$
B(Y') \longrightarrow B(Y) \times_{B(Z)} B(Z')
$$
is a bijection, see Pushouts of Spaces, Lemma
\ref{spaces-pushouts-lemma-pushout-along-thickening-schemes}.
Thus using the commutative diagram
$$
\xymatrix{
\Complexesstack_{X/B, Y'} \ar[r] \ar[d] &
\Complexesstack_{X/B, Y}
\times_{\Complexesstack_{X/B, Z}}
\Complexesstack_{X/B, Z'}
\ar[d] \\
B(Y') \ar[r] & B(Y) \times_{B(Z)} B(Z')
}
$$
we see that we may assume that $Y'$ is a scheme over $B'$. By
Remark \ref{remark-complexes-base-change}
we may replace $B$ by $Y'$ and $X$ by $X \times_B Y'$.
Thus we may assume $B = Y'$.

\medskip\noindent
Assume $B = Y'$. We first prove fully faithfulness of our functor.
To do this, let $\xi_1, \xi_2$ be two objects of $\Complexesstack_{X/B}$
over $Y'$. Then we have to show that
$$
\mathit{Isom}(\xi_1, \xi_2)(Y') \longrightarrow
\mathit{Isom}(\xi_1, \xi_2)(Y)
\times_{\mathit{Isom}(\xi_1, \xi_2)(Z)}
\mathit{Isom}(\xi_1, \xi_2)(Z')
$$
is bijective. However, we already know that $\mathit{Isom}(\xi_1, \xi_2)$
is an algebraic space over $B = Y'$. Thus this bijectivity follows from
Artin's Axioms, Lemma \ref{artin-lemma-pushout} (or the aforementioned
Pushouts of Spaces, Lemma
\ref{spaces-pushouts-lemma-pushout-along-thickening-schemes}).

\medskip\noindent
Essential surjectivity. Let $(E_Y, E_{Z'}, \alpha)$ be a triple,
where $E_Y \in D(\mathcal{O}_Y)$ and $E_{Z'} \in D(\mathcal{O}_{X_{Z'}})$
are objects such that $(Y, Y \to B, E_Y)$ is an object of
$\Complexesstack_{X/B}$ over $Y$, such that
$(Z', Z' \to B, E_{Z'})$ is an object of $\Complexesstack_{X/B}$ over $Z'$,
and $\alpha : L(X_Z \to X_Y)^*E_Y \to L(X_Z \to X_{Z'})^*E_{Z'}$
is an isomorphism in $D(\mathcal{O}_{Z'})$.
That is to say
$$
((Y, Y \to B, E_Y), (Z', Z' \to B, E_{Z'}), \alpha)
$$
is an object of the target of the arrow of our lemma.
Observe that the diagram
$$
\xymatrix{
X_Z \ar[r] \ar[d] & X_{Z'} \ar[d] \\
X_Y \ar[r] & X_{Y'}
}
$$
is a pushout with $X_Z \to X_Y$ affine and $X_Z \to X_{Z'}$ a thickening
(see Pushouts of Spaces, Lemma
\ref{spaces-pushouts-lemma-equivalence-categories-spaces-pushout-flat}).
Hence by Pushouts of Spaces, Lemma
\ref{spaces-pushouts-lemma-pushout-along-thickening-derived}
we find an object $E_{Y'} \in D(\mathcal{O}_{X_{Y'}})$
together with isomorphisms
$L(X_Y \to X_{Y'})^*E_{Y'} \to E_Y$ and
$L(X_{Z'} \to X_{Y'})^*E_{Y'} \to E_Z$
compatible with $\alpha$. Clearly, if we show that
$E_{Y'}$ is $Y'$-perfect, then we are done, because property (2)
of Lemma \ref{lemma-complexes}
is a property on points (and $Y$ and $Y'$ have the same points).
This follows from More on Morphisms of Spaces, Lemma
\ref{spaces-more-morphisms-lemma-thickening-relatively-perfect}.
\end{proof}

\end{document}
